\section{Projects}

\begin{apicard}{\method{GET}{getc}\quad\route{/projects}}
List all projects visible to the authenticated user, with pagination.

\vspace{3pt}\textbf{Query parameters}
\begin{tabularx}{\linewidth}{@{}l l X@{}}
\toprule
\textbf{Name} & \textbf{Type} & \textbf{Description}\\
\midrule
\ttfamily limit  & integer & Max items per page (1--100, default 20).\\
\ttfamily cursor & string  & Opaque pagination cursor from a previous response.\\
\ttfamily status & string  & Filter by \texttt{active}, \texttt{archived}, or \texttt{all}.\\
\bottomrule
\end{tabularx}

\vspace{4pt}\textbf{Response 200}
\begin{lstlisting}
{
  "data": [
    { "id": "prj_8a2f", "name": "Website Redesign",
      "status": "active", "created_at": "2026-06-01T09:12:00Z" }
  ],
  "next_cursor": "eyJvIjoyMH0"
}
\end{lstlisting}
\end{apicard}

\vspace{6pt}
\begin{apicard}{\method{POST}{postc}\quad\route{/projects}}
Create a new project.

\vspace{3pt}\textbf{Body parameters}
\begin{tabularx}{\linewidth}{@{}l l c X@{}}
\toprule
\textbf{Name} & \textbf{Type} & \textbf{Req.} & \textbf{Description}\\
\midrule
\ttfamily name        & string & yes & Human-readable project name.\\
\ttfamily description & string & no  & Optional summary.\\
\ttfamily template\_id & string & no  & Clone from an existing template.\\
\bottomrule
\end{tabularx}

\vspace{4pt}\textbf{Response 201}
\begin{lstlisting}
{ "id": "prj_9c7d", "name": "Mobile App",
  "status": "active", "created_at": "2026-06-28T14:03:00Z" }
\end{lstlisting}
\end{apicard}

% ================= ENDPOINT 2 =================
